\subsection{Polynomial dominance from a genuine central fibre}

\begin{lemma}[The actual zero-section point]\label{scheme:polynomial-zero-section-point-ideal}
\lean{EnvelopeProperty.PolynomialDominanceFromCentralFiber.zeroSectionPoint_asIdeal}
\leanok
Let \(k\) be a field. The image of the unique point of \(\operatorname{Spec}k\)
under the morphism induced by evaluation \(k[t]\to k,\ t\mapsto0\), is the prime ideal
\((t)\) of \(k[t]\).
\end{lemma}
\begin{proof}
The unique prime of the field \(k\) is zero, and the underlying map of spectra is
contraction. Its image ideal is therefore the kernel of evaluation at zero.
Division by the monic polynomial \(t\) shows that this kernel is exactly \((t)\).
These are the actual spectrum map and its actual prime ideal.
\end{proof}

\begin{theorem}[Dominance from a central point and a horizontal generic point]
\label{scheme:polynomial-dominance-central-point}
\lean{EnvelopeProperty.PolynomialDominanceFromCentralFiber.isDominant_of_zeroSectionPoint}
\leanok
\uses{scheme:polynomial-zero-section-point-ideal}
Let \(D\) be an integral scheme over an arbitrary field \(k\), and let
\(p:D\to\operatorname{Spec}k[t]\) be an actual scheme morphism. Suppose there is an
actual point \(z_0\in D\) mapping to the zero-section point, and that
\(t\notin p(\eta_D)\), where \(\eta_D\) is the generic point of \(D\).
Then \(p\) is dominant.
\end{theorem}
\begin{proof}
Set \(P=p(\eta_D)\), an actual prime ideal of \(k[t]\).
Every point of \(D\), including \(z_0\), is a specialization of \(\eta_D\).
Continuity of the actual morphism \(p\) therefore gives
\(p(\eta_D)\leadsto p(z_0)\). The specialization order on an affine spectrum is
inclusion of the corresponding prime ideals, so the preceding lemma gives
\(P\subseteq(t)\).

If \(P\ne0\), the principal ideal domain \(k[t]\) makes the nonzero prime \(P\)
maximal. The proper prime ideal \((t)\) containing it must then equal \(P\).
This would put \(t\) in \(P\), contradicting the given generic-point condition.
Consequently \(P=0\). Thus the actual image of the original generic point is
the generic point of \(\operatorname{Spec}k[t]\). The image of \(p\) contains
this point, whose closure is the whole target, proving dominance.

The central point is essential: the condition \(t\notin P\) alone does not
force \(P=0\). No closedness, finite-type, normality, flatness, projectivity,
or algebraic-closedness assumption is used.
\end{proof}

\begin{corollary}[Dominance from the whole original central-fibre triangle]
\label{scheme:polynomial-dominance-central-fibre}
\lean{EnvelopeProperty.PolynomialDominanceFromCentralFiber.isDominant_of_centralFiber}
\leanok
\uses{scheme:polynomial-dominance-central-point}
Let \(D\) be integral and \(C\) a nonempty scheme, over an arbitrary field \(k\).
Let actual scheme morphisms \(p:D\to\operatorname{Spec}k[t]\),
\(g:C\to D\), and \(q:C\to\operatorname{Spec}k\) satisfy the whole morphism identity
\[
 p\circ g=\operatorname{Spec}(\operatorname{ev}_0)\circ q
\]
from \(C\) to \(\operatorname{Spec}k[t]\).
If \(t\notin p(\eta_D)\), then \(p\) is dominant.
\end{corollary}
\begin{proof}
Choose a point \(c\in C\), using precisely the asserted nonemptiness.
Evaluate the given equality of the actual whole morphisms at \(c\).
The point \(q(c)\) is the unique point of \(\operatorname{Spec}k\), so
\(z_0=g(c)\) maps to the actual zero-section point. Apply the theorem to this
point and the unchanged original generic-point hypothesis.
No property of \(g\) such as closed immersion, and no irreducibility or
reducedness of \(C\), is needed.
\end{proof}

